%======================================
%卒業研究報告書サンプル
%ver. 1.0 2024-12-04 %このように履歴を残すのも良い
%ver. 1.1 2025-02-?? 投稿
%ver. 2.1 2025-02-?? 修正
%======================================

\documentclass[12pt,a4j,twoside]{jreport}

\usepackage{amsmath,amssymb}
\usepackage{amsthm}
\usepackage[abbrev]{amsrefs}
\usepackage{latexsym}


\pagestyle{plain}
\oddsidemargin 23pt
\evensidemargin 18pt
\renewcommand{\bibname}{参考文献}
\renewcommand{\proofname}{\bf 証明}
\renewcommand{\labelenumi}{{\rm(\arabic{enumi})}}
\numberwithin{equation}{chapter}

\theoremstyle{definition} % 以下の環境で英文がイタリックにならないようにするため. 
\newtheorem{thm}{定理}[chapter]
\newtheorem{prop}[thm]{命題}
\newtheorem{lem}[thm]{補題}
\newtheorem{cor}[thm]{系}
\newtheorem{dfn}[thm]{定義}
\newtheorem{ex}[thm]{例}
\newtheorem{rem}[thm]{注意}

%%newcommand
\newcommand{\R}{{\mathbb R}} %Euclid space
\newcommand{\Lap}{\Delta} %Laplacian
\newcommand{\FL}[2]{(-\Lap)^{#1\frac{#2}{2}}} %Fractional Laplacian
\newcommand{\Fl}[1]{(-\Lap)^{#1/2}} %Fractional Laplacian
\newcommand{\F}{\mathcal F} %Fourier transform
\newcommand{\gm}{\gamma} %gamma
\newcommand{\Gm}{\Gamma} %Gamma 
\newcommand{\conj}[1]{\overline{#1}} %conjugate
\newcommand{\Sp}{{\mathbb S}} %Sphere
\newcommand{\wb}{\qed} %Result end
\newcommand{\bb}{\hfill $\blacksquare$} %Proof end
\newcommand{\al}{\alpha} %alpha
\newcommand{\lm}{\lambda} %lambda
\renewcommand{\th}{\theta} %theta
\newcommand{\pt}{\partial} %partial derivative
\newcommand{\s}{\sigma} %sigma
\renewcommand{\t}{\tau} %tau
\newcommand{\Sw}{\mathcal S} %Schwartz space
\newcommand{\bs}{\backslash} %backslash
\newcommand{\Del}{\Delta} %Laplacian
\newcommand{\bt}{\beta} %beta
\newcommand{\ep}{\varepsilon} %epsilon
\newcommand{\C}{{\mathbb C}} %Complex number
\newcommand{\nab}{\nabla} %gradient
\newcommand{\dl}{\delta} %delta
\renewcommand{\r}{\rho} %rho
\newcommand{\N}{{\mathbb N}} %Natural number
\newcommand{\Om}{\Omega} %Omega
\newcommand{\kp}{\kappa} %kappa
\newcommand{\Hm}{\mathcal H} %Hausdorff measure
\newcommand{\om}{\omega} %omega
\newcommand{\ev}{\mathbf{e}_1}
\newcommand{\wc}{\rightharpoonup} %weak convergence
\newcommand{\D}{\mathcal D} %Distributional function space
\renewcommand{\max}{\mathop{\rm max}\limits} %max
\renewcommand{\sup}{\mathop{\rm sup}\limits} %sup
\renewcommand{\lim}{\mathop{\rm lim}\limits} %lim
\renewcommand{\div}{\mathrm{div}} %divergence
\newcommand{\diam}{\mathrm{diam}}
\renewcommand{\det}{\mathrm{det}} %deteminant
\newcommand{\rot}{\mathrm{rot}} %rotation
\newcommand{\Z}{\mathbb{Z}}
\newcommand{\He}{\mathbb{H}}
\renewcommand{\Im}{\mathrm{Im}}
\newcommand{\sL}{\mathcal{L}}

\DeclareMathOperator{\esssup}{ess, \sup}
\DeclareMathOperator{\real}{Re}
\DeclareMathOperator{\supp}{supp}
\DeclareMathOperator{\sgn}{sgn}
\DeclareMathOperator{\dist}{dist}
\DeclareMathOperator{\rank}{rank}

\renewcommand{\Re}{\real}

\newcommand{\eqn}[1]{\begin{equation}#1\end{equation}}
\newcommand{\eq}[1]{\begin{equation*}#1\end{equation*}}
\newcommand{\spln}[1]{\begin{equation}\begin{split}#1\end{split}\end{equation}}
\newcommand{\spl}[1]{\begin{equation*}\begin{split}#1\end{split}\end{equation*}}
\newcommand{\casn}[1]{\begin{eqnarray}\begin{cases}#1\end{cases}\end{eqnarray}}
\newcommand{\cas}[1]{\begin{eqnarray*}\begin{cases}#1\end{cases}\end{eqnarray*}}

\def\Xint#1{\mathchoice
{\XXint\displaystyle\textstyle{#1}}%
{\XXint\textstyle\scriptstyle{#1}}%
{\XXint\scriptstyle\scriptscriptstyle{#1}}%
{\XXint\scriptscriptstyle\scriptscriptstyle{#1}}%
\!\int}
\def\XXint#1#2#3{{\setbox0=\hbox{$#1{#2#3}{\int}$}
\vcenter{\hbox{$#2#3$}}\kern-.5\wd0}}
\def\ddashint{\Xint=}
\def\dashint{\Xint-}

\allowdisplaybreaks

\begin{document}

%======================================
\title{ %表紙
\huge 卒業研究の題目 \\[5cm]
}
\author{228s1341　江嶋 宏優 \\
\qquad\qquad\quad Ejima Atsuhiro\\[1cm]
（指導教員: 勝呂 剛志 准教授） \\[1cm]
熊本大学理学部理学科数学コース \\
卒業研究報告書}
\date{\Large 2027 年 2 月} %年・月を入れる. 
\maketitle %表紙を作成
\thispagestyle{empty} 
\clearpage
\tableofcontents %目次
\thispagestyle{empty}
\newpage
\setcounter{page}{1}
\renewcommand{\thepage}{\arabic{page}} 
%======================================


\chapter{序文}%ここに序文を書く. 

本稿では, 次の対数型Sobolevの不等式を考える: 

\begin{thm}[対数型Sobolevの不等式] \label{thm:dareka}
任意の$f \in H^1(\R^n)$に対して, 
\eqn{
   \int_{\R^n} |f(x)|^2 \log\frac{|f(x)|^2}{\|f\|_{L^2}^2}\, dx
   \leq \frac{n}{2} \|f\|_{L^2}^2 \log\left(\frac{2}{n \pi e \|f\|_{L^2}^2} \int_{\R^n} |\nab f(x)|^2\, dx\right) 
 }
\begin{enumerate}
\item これこれ. 
\item しかじか. 
\end{enumerate}
\end{thm}


以下で, $1 \leq p < \infty$に対して, $L^p$-空間を次で定める:
\begin{equation}\label{dfn;Lp}
  L^p(\Om)
  \equiv \left\{f: \Om \to \R;f: \text{可測},\quad \|f\|_{L^p} \equiv \left(\int_{\R^n} |f(x)|^p\, dx\right)^\frac{1}{p} < + \infty\right\}.
\end{equation}
%定義\eqref{dfn;Lp}より, $\|\cdot\|_{L^p}$はノルムなので, $L^p$はノルム空間である. 
$H^1$はSobolev空間とする. 


\begin{rem}
定数$2/(n \pi e)$は最良定数である.
最良定数はGauss函数によって達成されることが知られている. 
詳しくは, Carlen (1995)を参照せよ. 
\end{rem}


通常の$H^1$上では次のSobolevの不等式が成り立つことが知られている:


\begin{thm}
任意の$f \in H^1(\R^n)$に対して, 
\eqn{
  \|f\|_{L^\frac{2 n}{n - 2}}
  \leq C \|\nab f\|_{L^2}
}
が成り立つ. 
\end{thm}


Jensenの不等式より, 
\eqn{
  \int_{\R^n} |f(x)|^2 \log\frac{|f(x)|^2}{\|f\|_{L^2}^2}\, dx
  \leq C \|f\|_{L^\frac{2 n}{n - 2}}
}
が成り立つので, Sobolevの不等式から対数型Sobolevの不等式が導かれることがわかる. 


\begin{rem}
Sobolevの不等式の最良定数はTalenti, Liebによって同定されている. 
実は, Liebの証明においては, Hardy--Littlewood--Sobolevの不等式の最良定数の導出が鍵となる. 

\end{rem}


ガンマ函数は次で定まる:
\eqn{
  \Gm(x) = \int_0^\infty t^{x - 1} e^{- t}\, dt.
}
\eqn{
  \left|\Sp^{n - 1}\right|
  = \frac{2 \pi^\frac{n}{2}}{\Gm\left(\frac{n}{2}\right)}.
}




\begin{rem}
HLS不等式の最良定数は知られており, 
\end{rem}

定理\ref{thm:dareka}より以下が従う. 

\begin{cor}
系はこのように書く. 
\end{cor}

\begin{proof}
証明はこのように書く.（普通は序文では使わない. ）
\end{proof}

\section*{謝辞}
ここに謝辞を書く. 

\chapter{準備, Lebesgue空間}

ここに2章の内容を書く. 

\section{??}

\begin{prop}
命題はこのように書く. 
\end{prop}
\begin{proof}
証明はこのように書く. 証明終わりの目印などは
変更しても良いが全体で統一すること. 
\end{proof}
\begin{lem}
補題はこのように書く. 
\end{lem}
\begin{ex}
例はこのように書く. 
\end{ex}
\begin{dfn}
定義はこのように. 
\end{dfn}

\begin{rem}
注意や例, 定義も, 定理や補題などと同じ各章で通し番号となっている. 
各自で例えば全体で通し番号となるように変更しても良い. 
\end{rem}


\begin{align*}
f(x)
&= g(x) \\
&= g(x) \\
&= g(x) \\
&= g(x) \\&= g(x) \\
&= g(x) \\&= g(x) \\
&= g(x) \\&= g(x) \\
&= g(x) \\&= g(x) \\
&= g(x) \\&= g(x) \\
&= g(x) \\&= g(x) \\
&= g(x) \\
\end{align*}

\section{??}

文献の引用は \cite{Gmv:green}, \cite{APS} のように書く. 
%同一の著者による一連の文献をを引用する時は \cites{Satake, SatakeN} のように書いても良い. 

\chapter{主定理の証明}

ここに3章の内容を書く. 

\section{??}

\section{??}

１ページ分の分量はこのぐらいです. 
ご自身でサイズを調整したいときは, このサイズを
大幅に超えることのないようにしましょう. 
１ページ分の分量はこのぐらいです. 
ご自身でサイズを調整したいときは, このサイズを
大幅に超えることのないようにしましょう. 
１ページ分の分量はこのぐらいです. 
ご自身でサイズを調整したいときは, このサイズを
大幅に超えることのないようにしましょう. 
１ページ分の分量はこのぐらいです. 
ご自身でサイズを調整したいときは, このサイズを
大幅に超えることのないようにしましょう. 
１ページ分の分量はこのぐらいです. 
ご自身でサイズを調整したいときは, このサイズを
大幅に超えることのないようにしましょう. 
１ページ分の分量はこのぐらいです. 
ご自身でサイズを調整したいときは, このサイズを
大幅に超えることのないようにしましょう. 
１ページ分の分量はこのぐらいです. 
ご自身でサイズを調整したいときは, このサイズを
大幅に超えることのないようにしましょう. 
１ページ分の分量はこのぐらいです. 
ご自身でサイズを調整したいときは, このサイズを
大幅に超えることのないようにしましょう. 
１ページ分の分量はこのぐらいです. 
ご自身でサイズを調整したいときは, このサイズを
大幅に超えることのないようにしましょう. 
１ページ分の分量はこのぐらいです. 
ご自身でサイズを調整したいときは, このサイズを
大幅に超えることのないようにしましょう. 
１ページ分の分量はこのぐらいです. 
ご自身でサイズを調整したいときは, このサイズを
大幅に超えることのないようにしましょう. 
１ページ分の分量はこのぐらいです. 
ご自身でサイズを調整したいときは, このサイズを
大幅に超えることのないようにしましょう. 
１ページ分の分量はこのぐらいです. 
ご自身でサイズを調整したいときは, このサイズを
大幅に超えることのないようにしましょう. 
１ページ分の分量はこのぐらいです. 
ご自身でサイズを調整したいときは, このサイズを
大幅に超えることのないようにしましょう. 
１ページ分の分量はこのぐらいです. 
ご自身でサイズを調整したいときは, このサイズを
大幅に超えることのないようにしましょう. 




新しいページになりました. 

修士論文は推敲に推敲を重ねて十分検討して執筆しましょう. 
修士論文は推敲に推敲を重ねて十分検討して執筆しましょう. 
修士論文は推敲に推敲を重ねて十分検討して執筆しましょう. 
修士論文は推敲に推敲を重ねて十分検討して執筆しましょう. 
修士論文は推敲に推敲を重ねて十分検討して執筆しましょう. 
修士論文は推敲に推敲を重ねて十分検討して執筆しましょう. 
修士論文は推敲に推敲を重ねて十分検討して執筆しましょう. 
修士論文は推敲に推敲を重ねて十分検討して執筆しましょう. 
修士論文は推敲に推敲を重ねて十分検討して執筆しましょう. 
修士論文は推敲に推敲を重ねて十分検討して執筆しましょう. 
修士論文は推敲に推敲を重ねて十分検討して執筆しましょう. 
修士論文は推敲に推敲を重ねて十分検討して執筆しましょう. 
修士論文は推敲に推敲を重ねて十分検討して執筆しましょう. 
修士論文は推敲に推敲を重ねて十分検討して執筆しましょう. 
修士論文は推敲に推敲を重ねて十分検討して執筆しましょう. 
修士論文は推敲に推敲を重ねて十分検討して執筆しましょう. 
修士論文は推敲に推敲を重ねて十分検討して執筆しましょう. 
修士論文は推敲に推敲を重ねて十分検討して執筆しましょう. 
修士論文は推敲に推敲を重ねて十分検討して執筆しましょう. 
修士論文は推敲に推敲を重ねて十分検討して執筆しましょう. 
修士論文は推敲に推敲を重ねて十分検討して執筆しましょう. 
修士論文は推敲に推敲を重ねて十分検討して執筆しましょう. 
修士論文は推敲に推敲を重ねて十分検討して執筆しましょう. 
修士論文は推敲に推敲を重ねて十分検討して執筆しましょう. 
修士論文は推敲に推敲を重ねて十分検討して執筆しましょう. 
修士論文は推敲に推敲を重ねて十分検討して執筆しましょう. 
修士論文は推敲に推敲を重ねて十分検討して執筆しましょう. 
修士論文は推敲に推敲を重ねて十分検討して執筆しましょう. 
修士論文は推敲に推敲を重ねて十分検討して執筆しましょう. 
修士論文は推敲に推敲を重ねて十分検討して執筆しましょう. 
修士論文は推敲に推敲を重ねて十分検討して執筆しましょう. 
修士論文は推敲に推敲を重ねて十分検討して執筆しましょう. 
修士論文は推敲に推敲を重ねて十分検討して執筆しましょう. 
修士論文は推敲に推敲を重ねて十分検討して執筆しましょう. 
修士論文は推敲に推敲を重ねて十分検討して執筆しましょう. 
修士論文は推敲に推敲を重ねて十分検討して執筆しましょう. 
修士論文は推敲に推敲を重ねて十分検討して執筆しましょう. 
修士論文は推敲に推敲を重ねて十分検討して執筆しましょう. 
修士論文は推敲に推敲を重ねて十分検討して執筆しましょう. 
修士論文は推敲に推敲を重ねて十分検討して執筆しましょう. 
修士論文は推敲に推敲を重ねて十分検討して執筆しましょう. 
修士論文は推敲に推敲を重ねて十分検討して執筆しましょう. 
修士論文は推敲に推敲を重ねて十分検討して執筆しましょう. 
修士論文は推敲に推敲を重ねて十分検討して執筆しましょう. 
修士論文は推敲に推敲を重ねて十分検討して執筆しましょう. 
修士論文は推敲に推敲を重ねて十分検討して執筆しましょう. 
修士論文は推敲に推敲を重ねて十分検討して執筆しましょう. 
修士論文は推敲に推敲を重ねて十分検討して執筆しましょう. 
修士論文は推敲に推敲を重ねて十分検討して執筆しましょう. 



\begin{bibdiv}
\begin{biblist}

%「参考文献」には, 欧文図書, 欧文論文, 
%和書, プレプリントの掲載例を示してある. 
%その他プロシーディング等の文献については, 
%これらの例やその他の文献の「参考文献」の記述方法を
%参照しながら, ご自身で記載すること. 


% 以下は MathSciNet の Batch Download を AMSRefs に選んでチェックを付け, 
% Retrieve Marked ボタンを押すと出てくる. 

\bib{APS}{article}{
   author={Atiyah, M. F.},
   author={Patodi, V. K.},
   author={Singer, I. M.},
   title={Spectral asymmetry and Riemannian geometry. I},
   journal={Math. Proc. Cambridge Philos. Soc.},
   volume={77},
   date={1975},
   pages={43--69},
   issn={0305-0041},
%   review={\MR{0397797}},
}


\bib{AS}{article}{
   author={Atiyah, M. F.},
   author={Singer, I. M.},
   title={The index of elliptic operators. I},
   journal={Ann. of Math. (2)},
   volume={87},
   date={1968},
   pages={484--530},
   issn={0003-486X},
%   review={\MR{0236950}},
}


\bib{Gmv:green}{book}{
   author={Gromov, Misha},
   title={Metric structures for Riemannian and non-Riemannian spaces},
   series={Modern Birkh\"auser Classics},
   edition={Reprint of the 2001 English edition},
   note={Based on the 1981 French original;
   With appendices by M. Katz, P. Pansu and S. Semmes;
   Translated from the French by Sean Michael Bates},
   publisher={Birkh\"auser Boston, Inc., Boston, MA},
   date={2007},
   pages={xx+585},
   isbn={978-0-8176-4582-3},
   isbn={0-8176-4582-9},
%   review={\MR{2307192 (2007k:53049)}},
}

%和書は MathSciNet からは出力出来ない. 以下のように記述する. 

%\bib{Satake}{book}{
%  author={佐武一郎},
%   title={『線型代数学』},
%   series={数学選書 1},
%   publisher={裳華房},
%   date={1974年},
%}


%\bib{SatakeN}{book}{
%  author={佐武一郎},
%   title={『線型代数学』{\rm (新装版)}},
%   edition={(新装版)}, エディッションの記載は文献による
%   series={数学選書 1},
%   publisher={裳華房},
%   date={2015年},
%}

% プレプリントは以下のように記述
\bib{Preprint}{article}{
   author={Tao, T.},
   author={Ziegler, T.},
   title={Polynomial patterns in the primes},
   note={preprint (2016), arXiv:1603.07817},
}

\end{biblist}
\end{bibdiv}

\end{document}


\bye


\bye
